\section{Conclusion}
\label{sec:conclusion-final}
%=============================================================================

\subsection{Summary of Rigorous Results}

This paper has established the following \textbf{rigorous} results for four-dimensional $SU(N)$ lattice Yang--Mills theory:

\begin{theorem}[Rigorous Results---Finite Volume]
\label{thm:finite-vol-rigorous}
For any finite lattice $\Lambda$ of size $L$ and any coupling $\beta > 0$:
\begin{enumerate}
    \item The transfer matrix $T_L$ is a compact, positive operator with discrete spectrum.
    \item The largest eigenvalue $\lambda_0 = 1$ is simple (Perron-Frobenius).
    \item The spectral gap $\Delta_L(\beta) = -\log \lambda_1(L) > 0$.
    \item The theory satisfies reflection positivity.
\end{enumerate}
\end{theorem}

\begin{theorem}[Rigorous Results---Strong Coupling]
\label{thm:strong-coupling-rigorous-conclusion}
For $\beta < \beta_0 = c/N^2$ in the infinite-volume limit:
\begin{enumerate}
    \item The mass gap $\Delta(\beta) > 0$ (via cluster expansion).
    \item The string tension $\sigma(\beta) \geq -\log(I_1(\beta N)/I_0(\beta N)) > 0$.
    \item Correlation functions decay exponentially.
    \item The free energy $f(\beta)$ is real-analytic.
\end{enumerate}
\end{theorem}

\subsection{Summary of Open Problems}

The following problems remain \textbf{open}:

\begin{enumerate}
    \item \textbf{Mass gap for $\beta \geq \beta_0$:} Extending the spectral gap from strong coupling to all couplings.
    
    \item \textbf{Uniform-in-$L$ bounds:} Proving that $\Delta_L(\beta) \geq c(\beta) > 0$ uniformly in $L$ for fixed $\beta > \beta_0$.
    
    \item \textbf{Continuum limit existence:} Rigorous construction of the continuum Yang-Mills measure.
    
    \item \textbf{Continuum mass gap:} Proving $\Delta_{\text{phys}} > 0$ in the continuum theory.
    
    \item \textbf{Mass gap--string tension relation:} Rigorous derivation of $\Delta \geq c_N \sqrt{\sigma}$.
\end{enumerate}

\subsection{Conditional Results}

Under additional assumptions, we have shown:

\begin{theorem}[Conditional Mass Gap]
\label{thm:conditional-conclusion}
If:
\begin{enumerate}[label=(\alph*)]
    \item No phase transition occurs for any $\beta > 0$ (Conjecture~\ref{conj:no-phase-transition}), and
    \item The spectral gap is uniform in $L$ (Conjecture~\ref{conj:uniform-gap}),
\end{enumerate}
then the physical mass gap satisfies $\Delta_{\text{phys}} > 0$.
\end{theorem}

\subsection{Honest Assessment}

\begin{tcolorbox}[colback=red!5!white, colframe=red!75!black, title=\textbf{What This Paper Does Not Prove}]
\begin{enumerate}
    \item \textbf{The Millennium Prize problem is NOT solved.}
    \item The mass gap for $\beta > \beta_0$ is \textbf{not proved}.
    \item The continuum limit existence is \textbf{not proved}.
    \item The Giles-Teper relation is NOT rigorously derived.
\end{enumerate}
\end{tcolorbox}

\subsection{Mathematical Contributions}

Despite not solving the full problem, this paper makes the following contributions:

\begin{enumerate}
    \item \textbf{Explicit bounds at strong coupling:} Quantitative estimates with explicit constants.
    
    \item \textbf{Clear identification of obstacles:} Precise statements of what remains to be proven.
    
    \item \textbf{Separation of rigorous from conjectural:} Honest assessment of the status of each claim.
    
    \item \textbf{Framework for future work:} Mathematical structure that could support further progress.
\end{enumerate}

\subsection{Directions for Future Research}

Progress on the Yang-Mills mass gap problem may come from:

\begin{enumerate}
    \item \textbf{Balaban's renormalization group:} Completing the program to control all coupling regimes.
    
    \item \textbf{Stochastic quantization:} Using Langevin dynamics with uniform ergodicity estimates.
    
    \item \textbf{New variational methods:} Constructive bounds on the spectral gap from below.
    
    \item \textbf{Supersymmetric interpolation:} Rigorous version of the adjoint QCD argument.
    
    \item \textbf{Lattice improvements:} Actions designed to have better analytical properties.
\end{enumerate}

\subsection{Final Remarks}

The Yang-Mills existence and mass gap problem remains one of the outstanding challenges in mathematical physics. While we have made progress in understanding the strong coupling regime and identifying the precise obstacles, the full problem---proving the mass gap in the continuum limit---is still open.

We hope that the honest assessment provided in this paper will help focus future research on the key mathematical difficulties that must be overcome.

\vspace{1cm}

\begin{center}
\fbox{\parbox{0.8\textwidth}{
\textbf{Status:} The Millennium Prize problem for Yang-Mills theory remains \textbf{OPEN}.

\textbf{This paper:} Provides rigorous results at strong coupling and identifies open problems.
}}
\end{center}


